\documentclass[11pt]{article}
\usepackage{mathblatt}
\makeatletter
\setlength{\mbpfbe}{0mm}% keine Punkte (Bauregel 6.4)
\makeatother
% Rechenraum als Karo (Bauregel 6.9): n Rechenschritte = n mal 10 mm
\newcommand{\karo}[1]{\par\smallskip\noindent\begin{tikzpicture}
  \clip (0,0) rectangle ({\linewidth-1mm},{#1*10mm});
  \draw[black!16,line width=0.3pt,step=5mm] (0,0) grid ({\linewidth},{#1*10mm});
  \end{tikzpicture}\par}
\newcommand{\linie}[1]{\,{\color{black!45}\rule[-1pt]{#1}{0.4pt}}\,}
\newcommand{\pkt}{\textperiodcentered{}}
\begin{document}

\pfheftstil
\pflaufkopf[P10 \pkt{} mehr]{Gleichungssystem im Sachzusammenhang \pkt{} Lösungen}{PA20-1}\pfseitevon
\pfheftkopf{Gleichungssystem im Sachzusammenhang}{Lösungen \pkt{} mehr}
\begin{pfloesungb}{}
\lzb{1.}{}{I: 2x + 2y = 61,80; Erwachsener 24,60 €; Kind 12,40 €}{I: 2x + 2y = 74,00, II: x + 3y = 61,80; Erwachsener 24,60 €; Kind 12,40 €}
\end{pfloesungb}
\begin{pfloesungb}{}
\lzb{2.}{}{6 Drei-Bett-Zimmer; 8 Fünf-Bett-Zimmer}{6 Drei-Bett-Zimmer; 8 Fünf-Bett-Zimmer}
\end{pfloesungb}
\end{document}